\cchapter{ЗАКЛЮЧЕНИЕ}                       % Заголовок
\addcontentsline{toc}{chapter}{ЗАКЛЮЧЕНИЕ}  % Добавляем его в оглавление

Данная работа позволила достигнуть поставленной цели - повышения эффективности процесса оценки стоимости ремонта поврежденных автомобилей. Достижение данной цели оценивалось по двум основным факторам:
\begin{enumerate}[label=\arabic*.]
  \item Производительность системы
  \item Минимизация человеческого участия
\end{enumerate}

\begin{table}[h!]
\centering
\begin{tabular}{|p{5cm}|p{3cm}|p{5cm}|p{3cm}|}
\hline
\textbf{Критерий} & \textbf{Исходное состояние} & \textbf{Результат после внедрения} & \textbf{Вывод} \\
\hline
Скорость получения предварительной сметы и длительности работ & 30 минут & 1 минута & Значительно сократилась \\
\hline
Возможность параллельной обработки обращений & Ограниченная & Есть за счет микросервисной архитектуры и очередей & Цель достигнута \\
\hline
Надежность и воспроизводимость сценария & Зависела от сотрудника & Формализована в виде системы и тестов & Цель достигнута \\
\hline
\end{tabular}
\caption{Сравнение состояния до и после внедрения системы}
\end{table}

С точки зрения производительности разработанный сервис обеспечивает автоматизированное прохождение сценария обработки обращения. Полученные в нагрузочном тестировании значения по пропускной способности и времени отклика показывают, что система способна выдержать поток обращений на уровне не менее 120 RPS и оставаться отказоустойчивой. Следовательно, первый критерий достижения цели можно считать выполненным.

С точки зрения минимизации человеческого участия цель также достигнута. В разработанном сервисе пользователь получает предварительный результат автоматически, и сотруднику автосервиса не требуется тратить время и отвлекаться на согласования и переписку с клиентом. Скорость получения ответа на интересующие вопросы клиента сокращена с 30 минут до 1 минуты, что высвобождает время сотрудников. Это позволит существенно оптимизировать текущий бизнес-процесс.

Степень достижения цели можно оценить как высокую, но не предельную. Разработанная система имеет перспективы дальнейшего развития:

\begin{enumerate}[label=\arabic*.]
  \item Качество распознавания зависит от объёма и качества размеченных данных. Необходимо расширение выборки, особенно по редким и сложным классам.
  \item Текущая версия системы анализирует одно изображение в рамках одного сценария, тогда как в реальной практике для более точной оценки может потребоваться несколько фотографий с разных ракурсов. Развитие системы в сторону анализа нескольких фотографий позволит избежать такого рода проблем.
  \item В текущей версии системы проводится классификация повреждений без оценки тяжести дефектов. Например, оценка глубины или площади вмятин позволила бы точнее рассчитывать стоимость ремонта.
\end{enumerate}

% Заключение должно содержать:
% \begin{itemize}
%   \item краткие выводы по результатам выполненной ВКР или отдельных её этапов;
%   \item оценку полноты решений поставленных задач;
%   \item разработку рекомендаций и исходных данных по конкретному использованию результатов ВКР;
%   \item результаты оценки технико-экономической эффективности внедрения (если имеет место);
%   \item результаты оценки научно-технического уровня выполненной ВКР в сравнении с достижениями в этой области.
% \end{itemize}
